\chapter{Phonons II: Thermal Properties}

The thermal properties of solids---heat capacity, thermal conductivity, and thermal expansion---are governed by lattice vibrations. The experimental data presented in this chapter demonstrate how phonon statistics and scattering processes determine these properties across an enormous range of temperatures.

\section{Phonon Heat Capacity}

\subsection{The Debye Model and the $T^3$ Law}

At low temperatures, the heat capacity of a solid follows the \textbf{Debye $T^3$ law}:
\begin{equation}
  C_V = \frac{12\pi^4}{5} N k_B \left(\frac{T}{\Theta_D}\right)^3 = 234\, N k_B \left(\frac{T}{\Theta_D}\right)^3
\end{equation}
where $\Theta_D$ is the \textbf{Debye temperature}, a material-specific parameter related to the maximum phonon frequency: $k_B \Theta_D = \hbar\omega_D$.

\paragraph{Experimental verification.} The $T^3$ law has been verified for hundreds of materials. The standard test is to plot $C/T$ vs.\ $T^2$ at low temperatures---the data should fall on a straight line:
\begin{equation}
  \frac{C}{T} = \gamma + \beta T^2
\end{equation}
where $\gamma T$ is the electronic contribution (Chapter~6) and $\beta = (12\pi^4/5)(Nk_B/\Theta_D^3)$ gives the Debye temperature. Parkinson (1958) performed this analysis for many metals, showing that the Debye model fits the data remarkably well below $\Theta_D/10$.

\subsection{Debye Temperatures}

The Debye temperature is one of the most fundamental parameters characterizing a solid. It correlates with bond strength, melting point, and sound velocity.

\begin{table}[H]
\centering
\caption{Experimental Debye temperatures (\si{\kelvin}) for selected elements and compounds. Values from Abrahams et al.\ (1975) and Grimvall et al.\ (1974).}
\label{tab:debye}
\begin{tabular}{@{}lSlS@{}}
\toprule
Material & {$\Theta_D$ (\si{\kelvin})} & Material & {$\Theta_D$ (\si{\kelvin})} \\
\midrule
C (diamond)  & 2230 & Cu    & 343 \\
Si           & 645  & Au    & 165 \\
Ge           & 374  & Al    & 428 \\
SiC          & 1200 & Fe    & 470 \\
MgO          & 946  & W     & 400 \\
NaCl         & 321  & Na    & 158 \\
KCl          & 235  & Pb    & 105 \\
LiF          & 732  & Ar (solid) & 93 \\
\bottomrule
\end{tabular}
\end{table}

Diamond has the highest Debye temperature of any material ($\Theta_D = \SI{2230}{\kelvin}$), reflecting its extremely stiff covalent bonds and light carbon atoms. At room temperature ($T/\Theta_D \approx 0.13$), diamond is still far from the classical limit---this is why Einstein's 1907 quantum theory of specific heat was inspired by diamond's anomalously low heat capacity.

\subsection{Einstein Model}

Einstein (1907) proposed the first quantum theory of heat capacity, treating all atoms as independent oscillators at the same frequency $\omega_E$:
\begin{equation}
  C_V = 3Nk_B \left(\frac{\Theta_E}{T}\right)^2 \frac{e^{\Theta_E/T}}{(e^{\Theta_E/T} - 1)^2}
\end{equation}
where $\Theta_E = \hbar\omega_E/k_B$ is the Einstein temperature.

The Einstein model correctly predicts the exponential freezing out of heat capacity at $T \ll \Theta_E$, but gives $C \propto e^{-\Theta_E/T}$ instead of the correct $T^3$ behavior. The Debye model, which accounts for the distribution of phonon frequencies, corrects this.


\section{Thermal Conductivity}

\subsection{The Phonon Gas Model}

In a nonmetallic crystal, heat is carried by phonons. The thermal conductivity is:
\begin{equation}
  \kappa = \frac{1}{3} C_V v_s \ell
\end{equation}
where $v_s$ is the average phonon velocity and $\ell$ is the phonon mean free path.

The temperature dependence of $\kappa$ reflects the competition between three scattering mechanisms:
\begin{enumerate}[nosep]
  \item \textbf{Boundary scattering} (low $T$): $\ell \approx$ crystal size, $\kappa \propto T^3$ (following $C_V$).
  \item \textbf{Point defect scattering}: $\ell \propto \omega^{-4}$ (Rayleigh scattering).
  \item \textbf{Phonon--phonon (Umklapp) scattering} (high $T$): $\ell \propto e^{\Theta_D/bT}$, $\kappa$ decreases.
\end{enumerate}

The result is a characteristic \textbf{thermal conductivity peak} at a temperature typically $\Theta_D/20$ to $\Theta_D/10$.

\subsection{Diamond: The Highest Thermal Conductivity}

Diamond has the highest thermal conductivity of any known bulk material, reaching $\kappa > \SI{10000}{\watt\per\metre\per\kelvin}$ in isotopically pure crystals at the peak ($\sim$\SI{80}{\kelvin}).

Inyushkin et al.\ (2018) measured thermal conductivity of high-purity synthetic single-crystal diamonds from \SI{6}{\kelvin} to \SI{400}{\kelvin}, separating the contributions of normal and Umklapp phonon--phonon scattering:

\begin{table}[H]
\centering
\caption{Thermal conductivity of single-crystal diamond (\si{\watt\per\metre\per\kelvin}) at selected temperatures. Data from Inyushkin et al.\ (2018).}
\label{tab:diamond_kappa}
\begin{tabular}{@{}SS@{}}
\toprule
{$T$ (\si{\kelvin})} & {$\kappa$ (\si{\watt\per\metre\per\kelvin})} \\
\midrule
  6  &    15 \\
 10  &   120 \\
 20  &  2800 \\
 40  & 12000 \\
 80  & 14000 \\
150  &  4500 \\
300  &  2200 \\
400  &  1600 \\
\bottomrule
\end{tabular}
\end{table}

The enormous peak at \SI{80}{\kelvin} (where $\kappa$ reaches \SI{14000}{\watt\per\metre\per\kelvin}) and the subsequent rapid decrease demonstrate the onset of Umklapp scattering. At room temperature, $\kappa \approx \SI{2200}{\watt\per\metre\per\kelvin}$---still 5 times higher than copper.

\subsection{Umklapp Processes}

In a \textbf{normal} (N) process, two phonons scatter to produce a third with total crystal momentum conserved within the first Brillouin zone: $\mathbf{q}_1 + \mathbf{q}_2 = \mathbf{q}_3$.

In an \textbf{Umklapp} (U) process, the resulting wavevector falls outside the first zone and is folded back by a reciprocal lattice vector: $\mathbf{q}_1 + \mathbf{q}_2 = \mathbf{q}_3 + \mathbf{G}$.

Only Umklapp processes create thermal resistance (N-processes conserve phonon momentum and cannot equilibrate a heat current). The rate of U-processes decreases exponentially at low temperatures ($\propto e^{-\Theta_D/bT}$) because they require phonons with $q$ near the zone boundary, which are thermally activated.

\begin{table}[H]
\centering
\caption{Room-temperature thermal conductivity of selected solids (\si{\watt\per\metre\per\kelvin}).}
\label{tab:kappa_room}
\begin{tabular}{@{}lSl@{}}
\toprule
Material & {$\kappa$ (\si{\watt\per\metre\per\kelvin})} & Dominant heat carrier \\
\midrule
Diamond        & 2200  & phonons \\
Cu             &  401  & electrons \\
Al             &  237  & electrons \\
Si             &  148  & phonons \\
Fe             &   80  & electrons + phonons \\
GaAs           &   46  & phonons \\
NaCl           &    7  & phonons \\
Glass (SiO$_2$)&    1.4 & phonons (localized) \\
\bottomrule
\end{tabular}
\end{table}

\vspace{1cm}
\noindent\rule{\textwidth}{0.4pt}
\textit{Having established the lattice dynamics, we now turn to the electronic properties of solids, beginning with the free electron model.}
